\chapter{Criptografía Clásica}

\section{Introducción}
\begin{itemize}
  \item \textbf{Criptografía simétrica}: misma clave para cifrar y descifrar.
  \item \textbf{Criptografía asimétrica}: claves diferentes.
\end{itemize}

\section{Cifrado César}
Desplaza cada letra un número fijo de posiciones. Con desplazamiento 3:
\[
\text{HOLA} \to \text{KROD}.
\]
Solo tiene 25 desplazamientos: es vulnerable a fuerza bruta.

\section{Cifrado de Sustitución}
Reemplaza cada letra por otra según una tabla. El \textbf{análisis de frecuencia} permite romperlo.

\section{Cifrado de Vigenère}
Usa una palabra clave para desplazar cada letra de forma diferente. Con clave \texttt{CLAVE}:
\[
\text{HOLA} \to \text{JPRD}.
\]

\section{Implementación en Haskell}

\begin{lstlisting}[caption={Criptografía clásica en Haskell}]
import Data.Char

cesar :: Int -> String -> String
cesar n = map (cifrar n)
  where
    cifrar n c
      | isAlpha c = chr (ord 'A' + ((ord (toUpper c) - ord 'A' + n) `mod` 26))
      | otherwise = c

cesarFuerzaBruta :: String -> [String]
cesarFuerzaBruta texto = [cesar n texto | n <- [0..25]]

vigenere :: String -> String -> String
vigenere clave texto =
    zipWith cifrar (map toUpper texto) (cycle (map toUpper clave))
  where
    cifrar c k = chr (ord 'A' + ((ord c - ord 'A' + ord k - ord 'A') `mod` 26))
\end{lstlisting}

\section{Ejercicios}
\begin{enumerate}
  \item Cifrar y descifrar con César.
  \item Romper un cifrado César por fuerza bruta.
  \item Implementar Vigenère en Haskell.
\end{enumerate}